\documentclass[10pt,a4paper]{amsart}
\usepackage[margin=19mm]{geometry}
\usepackage{amsmath,amssymb,mathtools}
\usepackage[scr=rsfs,cal=euler]{mathalfa}
\usepackage{xfrac}
\usepackage[unicode,hidelinks,bookmarks=false]{hyperref}
\newcommand{\ie}{i.e.\ }
\newcommand{\cf}{cf.\ }
\newcommand{\resp}{respectively\ }
\newcommand{\Subsection}{\subsection}
\providecommand{\nobreakdash}{\nobreak\hbox{-}}
\setlength{\emergencystretch}{2em}
\multlinegap0pt
\usepackage{bm}
\usepackage{bbm}
\usepackage{dsfont}
\usepackage{xspace}
\usepackage{cancel}
\usepackage{mathtools}
\usepackage{amsthm}
\makeatletter
\@ifundefined{theorem}{\newtheorem{theorem}{Theorem}}{}
\@ifundefined{definition}{\newtheorem{definition}[theorem]{Definition}}{}
\@ifundefined{lemma}{\newtheorem{lemma}[theorem]{Lemma}}{}
\makeatother
\title[Two aspects of graph 3-manifold groups]{Two aspects of graph 3-manifold groups}
\author{Hongbin Sun}
\date{Source: arXiv:2607.04326v1 (CC BY 4.0)}
\begin{document}
\maketitle

\noindent\emph{Source and licence.} Two aspects of graph 3-manifold groups. Version arXiv:2607.04326v1 (CC BY 4.0), distributed under the Creative Commons Attribution 4.0 International licence, as recorded in the arXiv licence metadata for this exact version and shown on the abstract page https://arxiv.org/abs/2607.04326v1. Full rights notice: licenses/arxiv-2607.04326-CC-BY-4.0.txt.\par\smallskip

\noindent\textit{Editorial scope.} Counted unit: the Lemma whose source label is \texttt{starnproperty} in the article (source0.tex lines 1286--1293, source label \texttt{starnproperty}), delivered together with the article's own complete original proof of it (source0.tex lines 1295--1337, one \texttt{proof} environment of 43 lines). No mathematics is written, repaired or re-derived: statement and proof are the article's own text line for line. Required context retained uncounted: starn (lines 1265--1276). Every definition, display and label of those lines is retained unchanged. Omitted from this package and not redistributed: 1--1264, 1277--1285, 1294--1294, 1338--1511. Nothing inside a retained excerpt is omitted or altered: every retained body is delivered exactly as the article's lines are written, so no omission receipt of any kind accompanies this package. Citations: the retained excerpts carry no citation command at all, so no bibliography entry of the article is transcribed and no reference is redistributed. Reference adaptation: none was needed and none was made; every reference inside the delivered unit resolves inside it, so no unresolved reference or citation is produced. Printed slips kept literal: no printed word, symbol, hypothesis or number of the counted statement or of its proof was altered. Layout and class: the article's own class is \texttt{amsart} with packages including amssymb, amsmath, mathrsfs, amsfonts, amscd, amsthm, graphics, graphicx, mathabx, fancyhdr, xy, hyperref, multirow, enumitem; the wrapper is the lane's pinned \texttt{amsart} editorial wrapper at 10pt on A4, which restates the theorem-like environments the retained text needs and adds the article's own macro definitions that the retained text needs (none), with extra wrapper packages (bm, bbm, dsfont, xspace, cancel, mathtools). The delivered numbering is the unit's own. Assets: no retained span contains a figure environment, a graphics-inclusion command, a PDF-page inclusion, a table or a TikZ picture; no image file is copied into this package, the package declares no asset dependency and no image file is redistributed with it. Licence: version arXiv:2607.04326v1 of this article is distributed under the Creative Commons Attribution 4.0 International licence, as recorded in the arXiv licence metadata for this exact version and shown on the abstract page https://arxiv.org/abs/2607.04326v1. A full rights notice is delivered at licenses/arxiv-2607.04326-CC-BY-4.0.txt. Characters: every retained excerpt is pure ASCII, so the delivered engine, XeLaTeX with its default Latin Modern text font, reports no missing character.
\begin{definition}\label{starn}
For any integer $n\geq 0$, we say that a group $G$ satisfies {\it property $(\star_n)$} if it contains a sequence of subgroups 
\begin{align}\label{8.1}
G=G_0\vartriangleright  G_1\vartriangleright  G_2\vartriangleright  \cdots \vartriangleright  G_{2n}
\end{align}
such that the following hold.
\begin{enumerate}
\item For any $i=0,\cdots,n-1$, $G_{2i}/G_{2i+1}$ is a finite group.
\item For any $i=0,\cdots,n-1$, $G_{2i+1}/G_{2i+2}$ is an abelian group (possibly infinitely generated).
\item $G_{2n}$ is isomorphic to a free group (possibly infinitely generated).
\end{enumerate}
\end{definition}

\begin{lemma}\label{starnproperty}
Let $G,H$ be two groups satisfying property $(\star_n)$ for some integer $n\geq 0$.
\begin{enumerate}
\item If $A<G$ is a subgroup, then $A$ satisfies property $(\star_n)$.
\item If $G<B$ is a finite-index subgroup, then $B$ satisfies property $(\star_n)$.
\item $G*H$ satisfies property $(\star_n)$.
\end{enumerate}
\end{lemma}

\begin{proof}
We first prove item (1). Let 
$$G=G_0\vartriangleright G_1\vartriangleright G_2\vartriangleright \cdots \vartriangleright G_{2n}$$
be a sequence of subgroups satisfying Definition \ref{starn}. Then we take $A_i=A\cap G_i$ to get a sequence of subgroups
\begin{align}\label{8.2}
A=A_0>A_1>A_2>\cdots>A_{2n}.
\end{align}
Since $G_{i+1}\vartriangleleft  G_i$, it is clear that $A_{i+1}=A\cap G_{i+1}$ is a normal subgroup of $A_i=A\cap G_i$.
Moreover, we have
$$A_i/A_{i+1}=A\cap G_i/A\cap G_{i+1}\cong (A\cap G_i)\cdot G_{i+1}/G_{i+1}<G_i/G_{i+1}.$$
Since a subgroup of a finite/abelian group is still finite/abelian, $A_i/A_{i+1}$ satisfies the condition in Definition \ref{starn} (1) (2). Since $A_{2n}=A\cap G_{2n}$ is a subgroup of a free group $G_{2n}$, $A_{2n}$ is also a free group. So the sequence of subgroups in \eqref{8.2} satisfies the condition in Definition \ref{starn}, thus $A$ has property $(\star_n)$.

Now we prove item (2). Since $G_1<G$ and $G<B$ are both finite-index subgroups, there exists a finite-index normal subgroup $B_1\vartriangleleft B$ contained in $G_1$. Then we take $B_i=G_i\cap B_1$ for $i=2,\cdots,2n$, and we use 
$$B=B_0\vartriangleright B_1>B_2>\cdots>B_{2n}$$ 
to check that $B$ satisfies property $(\star_n)$, as in the proof of item (1). 

We prove item (3) by inductively proving the following stronger statement. For any integer $n\geq 0$, let $$G=*_{\lambda\in \Lambda}G_{\lambda}$$
be a free product such that $\{G_{\lambda}\}_{\lambda\in \Lambda}$ has only finitely many isomorphic types and each $G_{\lambda}$ satisfies property $(\star_n)$, then $G$ satisfies property $(\star_n)$. Here we have no restriction on the cardinality of $\Lambda$, but we will only apply the result to countable $\Lambda$, since we only care about finitely generated $3$-manifold groups.

If $n=0$, then definition of property $(\star_0)$ implies that each $G_{\lambda}$ is isomorphic to a free group, then $G=*_{\lambda\in \Lambda}G_{\lambda}$ is also isomorphic to a free group and satisfies property $(\star_0)$.

Now, suppose the above statement holds for $n$, and we want to prove it for $n+1$. Suppose that $G=*_{\lambda\in \Lambda}G_{\lambda}$ such that each $G_{\lambda}$ is isomorphic to one of $G_{\lambda_1},\cdots,G_{\lambda_k}$, and each $G_{\lambda_i}$ satisfies property $(\star_{n+1})$. For any $i=1,\cdots,k$, $G_{\lambda_i}$ has a sequence of subgroups satisfying Definition \ref{starn}:
$$G_{\lambda_i}=G_{\lambda_i,0}\vartriangleright G_{\lambda_i,1}\vartriangleright G_{\lambda_i,2}\vartriangleright \cdots \vartriangleright G_{\lambda_i,2n+2}.$$

We take the surjective homomorphism 
$$\phi:G=*_{\lambda\in \Lambda}G_{\lambda}\to \Pi_{i=1}^k G_{\lambda_i}/G_{\lambda_i,1}$$
that restricts to each free factor $G_{\lambda}$ as 
$$G_{\lambda}\to G_{\lambda_i}\to G_{\lambda_i}/G_{\lambda_i,1}\to \Pi_{i=1}^k G_{\lambda_i}/G_{\lambda_i,1}.$$
Here, the first homomorphism is an isomorphism to some $G_{\lambda_i}$ with $i=1,\cdots,k$, the second homomorphism is the quotient homomorphism, and the third one is the natural inclusion. Then we take $G_1=\text{ker}(\phi)$, which is a finite-index normal subgroup of $G$ since $\Pi_{i=1}^k G_{\lambda_i}/G_{\lambda_i,1}$ is a finite group. 

The Kurosh Subgroup Theorem implies that 
$$G_1\cong (*_{\lambda'\in \Lambda_1}G_{\lambda',1})*F$$ is a free product, where each $G_{\lambda',1}$ is isomorphic to one of $G_{\lambda_1,1},\cdots, G_{\lambda_k,1}$ and $F$ is a free group.
Then we take the surjective homomorphism 
$$\psi:G_1= (*_{\lambda'\in \Lambda_1}G_{\lambda',1})*F\to \Pi_{i=1}^k G_{\lambda_i,1}/G_{\lambda_i,2}$$
that restricts to $F$ trivially and restricts to each free factor $G_{\lambda',1}$ as 
$$G_{\lambda',1}\to G_{\lambda_i,1}\to G_{\lambda_i,1}/G_{\lambda_i,2}\to \Pi_{i=1}^k G_{\lambda_i,1}/G_{\lambda_i,2}.$$
Here, the first homomorphism is an isomorphism to some $G_{\lambda_i,1}$ with $i=1,\cdots,k$, the second and third are quotient and inclusion homomorphisms, respectively. Then we take $G_2=\text{ker}(\psi)\vartriangleleft G_1$, and $G_1/G_2\cong\Pi_{i=1}^k G_{\lambda_i,1}/G_{\lambda_i,2}$ is abelian. 

By the Kurosh Subgroup Theorem again, we have $$G_2\cong (*_{\lambda''\in \Lambda_2}G_{\lambda'',2})*F'$$
where each $G_{\lambda'',2}$ is isomorphic to one of $G_{\lambda_1,2},\cdots, G_{\lambda_k,2}$ and $F'$ is a free group. 

Since $G_{\lambda_1,2},\cdots, G_{\lambda_k,2}$ and $F'$ all satisfy property $(\star_n)$, the inductive hypothesis implies that $G_2$ satisfies property $(\star_n)$. So $G$ satisfies property $(\star_{n+1})$ and the induction is done.
\end{proof}

\end{document}
